%% capitulos/conclusao.tex

\chapter{Conclusão}
\label{ch:conclusao}

Este trabalho descreveu como a complexidade estatística dos pesos da camada
densa de redes neurais profundas se comporta durante o treinamento, no contexto
de classificação de imagens médicas com a arquitetura ResNet-18 e o conjunto
MedNIST. A abordagem articulou duas medidas complementares --- a complexidade
LMC \cite{lmc1995} e a entropia amostral SampEn \cite{richman2000} ---
calculadas a cada época sobre a camada densa, em execuções de controle, com
subamostragem do conjunto de treino e com ruído branco e rosa nas imagens de
treino.

\section{Principais achados}
\label{sec:conc-achados}

\pendente{síntese do comportamento da LMC e da SampEn em cada regime, a partir
do Capítulo~\ref{ch:resultados}.}

\section{Trabalhos Futuros}
\label{sec:futuros}

A partir das limitações identificadas, perfilam-se os desdobramentos abaixo.

\begin{enumerate}
  \item \textbf{Repositório de maior dificuldade.} Repetir os experimentos em
  um conjunto de dados menos saturável que o MedNIST, no qual a perda de
  generalização seja mais pronunciada, e comparar sobre a camada densa os mesmos
  regimes de treinamento.

  \item \textbf{Ordem de achatamento.} A SampEn$_{1D}$ depende da ordem em que
  a matriz da camada densa é achatada (Seção~\ref{sec:camada-densa}). Convém
  comparar a ordem por neurônio, adotada como padrão, com a ordem por entrada.

  \item \textbf{Análise multiescala} (MSE) \cite{costa2002,costa2005}
  dos pesos da camada densa, para investigar dependências em diferentes
  escalas.

  \item \textbf{Comparação com outras arquiteturas} (EfficientNet,
  Wide ResNet) para verificar se o comportamento observado na ResNet-18
  se mantém.

  \item \textbf{Análise de sensibilidade à regularização explícita}
  (\emph{weight decay}, \emph{dropout}, \emph{label smoothing}) para
  testar se a manipulação dirigida da propensão ao sobreajuste altera o
  comportamento das medidas na camada densa.

  \item \textbf{Dados substitutos (\emph{surrogate data})} como teste de
  hipótese nula sobre a trajetória de complexidade, permitindo
  caracterizar o quanto da estrutura observada se afasta do
  comportamento esperado sob inicialização aleatória mais ruído.

  \item \textbf{Uso prático.} Avaliar, a partir do comportamento descrito,
  se as medidas podem compor um critério auxiliar de parada em pipelines de
  treinamento, comparando custo computacional e desempenho com os critérios
  tradicionais baseados na perda de validação.
\end{enumerate}
